\documentclass[11pt, a4paper, twoside]{book}
\usepackage{fontspec}

\setmainfont{Times New Roman}       % font chữ chính (serif)
\setsansfont{Arial}                 % font không chân (sans-serif)
\setmonofont{Courier New}           % font mono (code, terminal)

% Gói hỗ trợ tiếng Việt
\usepackage[vietnamese]{babel}


% ==============================================================================
% 2. CÁC GÓI TOÁN HỌC VÀ KHOA HỌC
% ==============================================================================
\usepackage{amsmath, amssymb, amsfonts, amsthm, caption, tabularx, csquotes, enumitem}
\setlistdepth{9}
% ==============================================================================
% 3. BỐ CỤC VÀ ĐỊNH DẠNG TRANG
% ==============================================================================
\usepackage[
a4paper,
top=2.5cm,
bottom=2.5cm,
left=3cm,
right=2.5cm,
headsep=1cm,
footskip=1cm,
twoside
]{geometry}
\usepackage{enumitem}
\setlist{nosep}
\usepackage{graphicx}
\usepackage{float}
\graphicspath{{images/}}
\usepackage[dvipsnames]{xcolor} 
\usepackage{hyperref}
\hypersetup{
	colorlinks=true,
	linkcolor=MidnightBlue,
	filecolor=magenta,      
	urlcolor=RoyalBlue,
	citecolor=ForestGreen,
	pdftitle={Giáo trình NLP Toàn tập},
	pdfauthor={Tên tác giả},
}
\usepackage{booktabs}
\usepackage{longtable}
% ==============================================================================
% 4. TÙY BIẾN TIÊU ĐỀ
% ==============================================================================
\usepackage{fancyhdr}
\pagestyle{fancy}
\fancyhf{}
\fancyhead[L]{\leftmark}
\fancyfoot[C]{Công Thái Học (ngu)}
\usepackage{titlesec}
\titleformat{\part}[display]
{\normalfont\Huge\bfseries\centering}
{\partname~\thepart}
{20pt}
{\Huge}
\titlespacing*{\part}{0pt}{50pt}{40pt}

\titleformat{\chapter}[display]
{\normalfont\huge\bfseries}
{\chaptertitlename~\thechapter}
{20pt}
{\huge\bfseries}
\titlespacing*{\chapter}{0pt}{40pt}{30pt}

\titleformat{\section}
{\normalfont\Large\bfseries\color{MidnightBlue}}
{\thesection.}
{1em}
{}
\titleformat{\subsection}
{\normalfont\large\bfseries\color{OliveGreen}}
{\thesubsection.}
{1em}
{}
\titleformat{\subsubsection}
{\normalfont\normalsize\bfseries\color{Brown}}
{\thesubsubsection.}
{1em}
{}

% ==============================================================================
% 5. CODE & BOX
% ==============================================================================
\usepackage{minted}
\setminted{
	frame=lines,
	framesep=2mm,
	numbersep=0pt,
	baselinestretch=1.2,
	fontsize=\small,
	linenos,
	breaklines,
	breakanywhere
}
\usepackage[most]{tcolorbox}
\tcbuselibrary{skins}

\newtcbtheorem{definition}{Định nghĩa}{
	colback=blue!5!white,
	colframe=blue!75!black,
	fonttitle=\bfseries,
	breakable
}{def}
\newtcbtheorem{example}{Ví dụ}{
	colback=green!5!white,
	colframe=green!50!black,
	fonttitle=\bfseries,
	before skip=0.5em,   
	after skip=0.5em,   
	boxsep=1mm,          
	breakable
}{ex}
\newtcbtheorem[
auto counter, 
number within=chapter,
list inside=rcp
]{recipe}{Recipe \thetcbcounter}{
	colback=orange!5!white,
	colframe=orange!75!black,
	fonttitle=\bfseries,
	title={Recipe \thechapter.\arabic{tcbcounter}:}
}{rcp}

% ==============================================================================
% 6. LỆNH PHỤ
% ==============================================================================
\newcommand{\newtag}{\texorpdfstring{\textcolor{red}{\textnormal{\bfseries\small [MỚI \& SỬA ĐỔI]}}}{}}
\newcommand{\reftag}{\textcolor{blue}{\textnormal{\bfseries\small [THAM KHẢO]}}}

% Chèn hình minh họa kèm mô tả nội dung.
% Cú pháp: \bookimage[độ rộng]{mã-hình}{mô tả chi tiết nội dung hình cần có}
% - Nếu images/<mã-hình>.png hoặc .jpg tồn tại: chèn ảnh.
% - Nếu chưa có: hiện khung chờ kèm mô tả để dễ bổ sung ảnh sau.
\newcommand{\bookimage}[3][0.8\textwidth]{%
	\IfFileExists{images/#2.png}{\includegraphics[width=#1]{#2.png}}{%
		\IfFileExists{images/#2.jpg}{\includegraphics[width=#1]{#2.jpg}}{%
			\fcolorbox{orange!70!black}{orange!5}{\parbox{\dimexpr#1-2\fboxsep-2\fboxrule\relax}{\centering\small
					\textbf{[Hình đang chờ bổ sung: \texttt{\detokenize{#2}}]}\\[2pt]\textit{#3}}}}}}

% ==============================================================================
% 7. THÔNG TIN SÁCH
% ==============================================================================
\title{GIÁO TRÌNH NLP TOÀN TẬP \\ \large}
\author{Đinh Công Thái}
\date{\today}

%%%%%%%%%%%%%%%%%%%%%%%%%%%%%%%%%%%%%%%%%%%%%%%%%%%%%%%%%%%%%%%%%%%%%%%%%%%%%%%%